% theory_6d_hcs_chiral_qg.tex -- Frontier Working Note
% 6d Holomorphic Theory, E_3 Chiral Structure, and Quantum Groups
% Status: OUTLINE with mathematical content per section
% Date: 2026-04-12 (updated with session results)

\documentclass[11pt]{article}
\usepackage[localtheorems]{raeez-math-template}
\newcommand{\Eone}{E_1}
\newcommand{\Etwo}{E_2}
\newcommand{\Einf}{E_\infty}
\newcommand{\fgl}{\mathfrak{gl}}
\newcommand{\frakg}{\mathfrak{g}}
\newcommand{\cC}{\mathcal{C}}
\newcommand{\cF}{\mathcal{F}}
\newcommand{\cZ}{\mathcal{Z}}
\newcommand{\Rep}{\mathrm{Rep}}
\newcommand{\HH}{\mathrm{HH}}
\newcommand{\CoHA}{\mathrm{CoHA}}
\newcommand{\C}{\mathbb{C}}
\newcommand{\Z}{\mathbb{Z}}
\newcommand{\R}{\mathbb{R}}
\theoremstyle{plain}
\newtheorem{theorem}{Theorem}[section]
\newtheorem{conjecture}[theorem]{Conjecture}
\newtheorem{proposition}[theorem]{Proposition}
\theoremstyle{definition}
\newtheorem{construction}[theorem]{Construction}
\theoremstyle{remark}
\newtheorem{remark}[theorem]{Remark}

\title{6d Holomorphic Theory, $E_3$ Chiral Structure, and Quantum Groups:\\
A Frontier Working Note for Volume III}
\author{Working note}
\date{April 2026}

\begin{document}
\maketitle

\tableofcontents

%% ====================================================================
%% PART I: THE ESTABLISHED CORE
%% ====================================================================

\section{The Established Core}

\subsection{3d holomorphic CS $\to$ Kac--Moody (Costello--Gwilliam)}
% PROVED. Boundary of hol CS on C x R is V_k(g).
% E_inf-chiral (BD sense), E_1-ordered refinement (Vol II Swiss-cheese).
% Prop \ref{prop:hcs-as-ce-cochains}: CE cochains = bulk observables.
% At critical level: Feigin--Frenkel center = Fun(Op_{G^L}(D)).

\subsection{5d holomorphic CS $\to$ affine Yangian (Costello 2013)}
% PROVED (Thm \ref{thm:5d-boundary-yangian}).
% Omega-background (h_1,h_2,h_3), h_1+h_2+h_3=0.
% Structure function g(u) = prod(u-h_i)/prod(u+h_i).
% E_2-chiral factorization on C^2; E_1-chiral projection to curve.
% Y^+(gl_hat_1) = CoHA, the associative positive half.
% Full Yangian via Drinfeld double.

\subsection{Proch\'azka--Rap\v{c}\'ak: $Y(\widehat{\fgl}_1) \simeq W_{1+\infty}$}
% PROVED ELSEWHERE. The isomorphism identifies the affine Yangian
% with W_{1+inf} at c=1. Character: vacuum module contributes P(q).
% Drinfeld double character: M(q)^2 * P(q) (Thm \ref{thm:drinfeld-center-coha}(iii)).

\subsection{Drinfeld center passage (297 tests across 4 suites)}
% PROVED HERE (Thm \ref{thm:drinfeld-center-coha}).
% Z(Rep^{E_1}(Y^+)) = Rep^{E_2}(Y(gl_hat_1)).
% E_1 -> E_2 via center; the braiding comes from universal R-matrix.
% Content: half-braiding on charge-1,2 Fock modules; 8x8 YBE symbolic;
% unitarity R_{12}(u)R_{21}(-u)=1; classical r-matrix r(u)=-sigma_2/u.
% 297 tests: drinfeld_center (43), _yangian (84), _hocolim (76), _e1_cy3 (94).

\subsection{MO $R$-matrix comparison (72 tests)}
% PROVED HERE (Thm \ref{thm:mo-chiral-rmatrix}).
% 7 verification paths: diagonal eigenvalues prod_{s in lambda} g(u+c(s)),
% Hilb^2(C^2) check, unitarity, YBE, crossing symmetry,
% classical r-matrix, SV specialization.
% Scope: C^3 (toric), K3 x E (via Hilb^n(K3)).

\subsection{$E_1$-chiral bialgebra axioms (Section 7 of e1\_chiral\_algebras.tex)}
% PROVED HERE. The five-part (H1)--(H5) axiom system for an
% E_1-chiral Hopf algebra (A, mu, Delta_z, epsilon, eta, S):
%   (H1) E_1-chiral algebra: augmented algebra in M^{E_1}_C.
%   (H2) E_1-chiral coalgebra: z-parametrized coproduct Delta_z
%        with deconcatenation; counit epsilon, coaugmentation eta.
%   (H3) Bialgebra compatibility: Delta_z is an E_1-chiral algebra map.
%   (H4) Spectral coassociativity:
%        (Delta_w x id) . Delta_{z+w} = (id x Delta_z) . Delta_w.
%   (H5) Hopf axiom: mu . (S x id) . Delta_0 = eta . epsilon
%        = mu . (id x S) . Delta_0.
% All 5 axioms verified: (H4) encodes the YBE; (H5) determines S
% uniquely on connected algebras. The averaging functor
% av: E_1-ChirHopf -> E_inf-ChirCoalg forgets all Hopf structure:
% R-matrix, antipode, z-dependence destroyed by symmetrization.
% This proves Li's vertex bialgebras cannot carry quantum group structure.

\subsection{Chiral coproduct spin-2 (23 tests)}
% SCAFFOLDED. Test suite for the chiral coproduct
% Delta_ch: A^! -> A^! hat-otimes A^! at spin 2.
% Deviation from primitive by Dimofte K-matrix K_A(z).
% K-matrix diagnostic: G/L/C/M classification via K complexity.

\subsection{$E_n$ hierarchy from holomorphic CS (47+141 tests)}
% holomorphic_cs_chiral_engine.py (47 tests): full pipeline
%   hol CS on C^n -> boundary A -> B(A) -> A^! -> Rep^{E_2}(A^!)
% higher_cy_en_tower.py (141 tests): E_1 stabilization theorem,
%   Bott periodicity obstruction for d=1,...,16,
%   d=4 Pontryagin obstruction, d=5 Z_2 Sp-refinement.

\subsection{Universal bar cohomology formula $(1+t)^{3g}$ for classes L, C}
% PROVED HERE (holomorphic_cs_chiral_engine.py, test_e3_bar_betagamma.py).
% For E_3 bar complex of chiral algebras of class >= L with g generators
% per direction:
%   Class G (Heisenberg, g=0 effective): H* = P(q)^3 (infinite, formal).
%   Class L (Yangian, g=1):    H* = (1+t)^3 = H*(T^3),  dim 8  = 2^3.
%   Class C (betagamma, g=2):  H* = (1+t)^6 = H*(T^6),  dim 64 = 2^6.
% Pattern: H*(B_{E_3}(A)) = (1+t)^{3g} = H*(T^{3g}), dim = 2^{3g}.
% The three E_3 differentials d_1, d_2, d_3 each replace one direction
% with its Koszul dual; the spectral sequence mechanism is uniform
% across classes L and C.

\subsection{K3 Yangian degree-24 structure function (68 tests)}
% PROVED HERE (k3_yangian.py, test_k3_yangian.py: 68 tests).
% The K3 Yangian Y(g_{K3}) has structure function of degree 24:
%   g_{K3}(z) = prod_{i=1}^{24} (z - h_i) / (z + h_i)
% with 24 Mukai lattice parameters h_i satisfying:
%   CY_2 constraint: sum h_i = 0 (analogous to h_1+h_2+h_3=0 for C^3).
%   Signature (4,20) from Mukai lattice Lambda_Muk = U^3 + E_8(-1)^2.
% Explicit Kummer K3 parametrization: h_i = epsilon (i=1..4),
%   h_i = -epsilon/5 (i=5..24), giving Mukai norm 16/5.
% Verified: unitarity g(z)g(-z) = 1, diagonal R-matrix R_{ii}(z) =
%   (z-h_i)/(z+h_i), classical limit r_{ii} = -2h_i,
%   Newton power sum / elementary symmetric function cross-checks.
% Degree 24 (not 3 as for C^3) reflects the full Mukai lattice rank.

\subsection{The Kummer route explicit computation (85 tests)}
% PROVED HERE (swiss_cheese_chart_gluing.py: 85 tests).
% Conj \ref{conj:kummer-route}: equivariant factorization homology
% on T^4/Z_2 computes the K3 chiral algebra in five steps:
%   (i) Heisenberg on T^4: rank-4 Heisenberg H_4, 8 generators.
%   (ii) Z_2 orbifold: H_4^{Z_2} + 16 twisted sectors from fixed points.
%   (iii) Singular Kummer orbifold algebra A^{orb}, character =
%         1/prod(1-q^n)^8 + 16 q^{1/2} / prod(1-q^n)^2.
%         kappa_ch(A^{orb}) = 2 (already matching smooth K3).
%   (iv) A_1 resolution of 16 singularities: 40 - 16 = 24 = dim H*(K3,C).
%   (v) Resolved Kummer chiral algebra A_{K3}^{Kum} = H_8 + 16 R_i / gluing.
% Verified alongside center-hocolim obstruction tests (76 tests in
% drinfeld_center_hocolim.py).

\subsection{Center-hocolim explicit obstruction (76+85 tests)}
% PROVED HERE (Prop \ref{prop:center-hocolim}).
% The Drinfeld center fails to commute with hocolim. For CY_3
% with multi-chart atlas {(Q_alpha, W_alpha)}:
%   Z(hocolim_alpha Rep^{E_1}(CoHA_alpha))
%     =/= hocolim_alpha Z(Rep^{E_1}(CoHA_alpha)).
% Explicit obstruction values:
%   C^3: 0 (single chart, no gluing needed).
%   Conifold: 2 (global center has dim 1, hocolim has dim 3).
%   Local P^2: nonzero.
%   K3 x E: massive (controlled by full BKM superalgebra).
% Consequence: the E_2-braided category is a GLOBAL construction,
% not assembable chart-by-chart. 76 tests in drinfeld_center_hocolim.py
% and 85 tests in swiss_cheese_chart_gluing.py.

\subsection{MO stable envelopes bypass center-hocolim}
% The Maulik--Okounkov stable-envelope R-matrix (Thm \ref{thm:mo-chiral-rmatrix},
% Thm \ref{thm:k3e-mo-rmatrix}) provides an independent route to E_2 braiding
% that BYPASSES the center-hocolim obstruction entirely.
% Mechanism: MO constructs R(u) directly from stable envelopes on
% Hilb^n(X), requiring only the equivariant geometry of the Nakajima
% variety, rather than the Drinfeld center of local CoHA charts.
% For toric CY_3: MO R-matrix matches chiral R-matrix (7 verification paths,
% 104 tests in mo_rmatrix_k3e.py).
% For K3 x E: MO stable envelopes on Hilb^n(K3) give the R-matrix directly,
% avoiding the multi-chart hocolim entirely. The K3 Yangian with degree-24
% structure function is the output.
% Net: the center-hocolim obstruction blocks chart-by-chart E_2 assembly,
% but MO provides a single-step geometric E_2 structure.

%% ====================================================================
%% PART II: THE CONJECTURAL PROGRAMME
%% ====================================================================

\section{The Conjectural Programme}

\subsection{6d holomorphic theory $\to$ $E_3$ factorization}
% Conj \ref{conj:topological-e3-comparison}. Non-Lagrangian 6d theory on C^3.
% This is not standard hol CS (Rem \ref{rem:6d-not-lagrangian}): kinetic term
% Omega wedge A wedge dbar A is (3,2)-form, top on R^5 not R^6.
% Route: (a) hol twist of 6d (2,0) SCFT, (b) Costello-Li perturbative,
% (c) CFG factorization homology of E_3 local observables.
% Two E_3 sources (AP154): algebraic (Deligne on HH^*(A,A) for E_inf A)
% vs topological (framed little 3-disks on Conf_n(C^3)).
% Agree rationally under Kontsevich formality; differ integrally.

\subsection{Quantum toroidal from $E_3$ projection}
% Conj \ref{conj:6d-boundary-toroidal}. E_3 on C^3 projected to E_1 on curve
% gives U_{q,t}(gl_hat_hat_1) with q=e^{2pi i h_1/h_3}, t=e^{-2pi i h_2/h_3}.
% Second affinization (second hat) from E_3; first from E_2.
% Intermediate E_2-projection to C^2 = affine Yangian (Thm 5d).

\subsection{Borcherds lift = resummation of perturbative series}
% ESTABLISHED (k3_times_e.tex Sec perturbative factorization homology).
% The perturbative factorization homology int_{K3} F is computed
% order-by-order in sigma_3 (the string coupling expansion):
%   Tree level (sigma_3 = 0): lattice VOA V_{Lambda_{K3}}, class G.
%   One loop (sigma_3^1): cup product correction, class L, depth 3.
%   Two loops (sigma_3^2): second FJ coefficient phi_2 of Delta_5.
%   All loops (sigma_3^n -> infty): RESUMS to the Borcherds product
%     for Delta_5, producing class M (infinite shadow depth).
% The Borcherds multiplicative lift is the non-perturbative resummation
% of the perturbative chiral bar complex expansion.
% K3 formality (m_k = 0 for k >= 3) means the entire series is
% controlled by the single operation m_2 (cup product) at all loop orders.
% 42 tests in k3e_coha_structure.py (perturbative FH subsuite).

\subsection{Class M = mock modular forms}
% ESTABLISHED (k3_times_e.tex, working_notes.tex).
% Shadow depth class M (r_{max} = infty) for K3 x E corresponds to:
%   (i) BKM root multiplicities c(D) != 0 for infinitely many D,
%       with |c(D)| ~ e^{4 pi sqrt(D)} (Rademacher asymptotics).
%   (ii) The 1/4-BPS counting function h(tau) is a mock modular form
%        of weight 1/2, shadow = 24 eta^3 (chi(K3) * eta^3).
%   (iii) Mock modularity = wall-crossing: departure from modularity
%         measures chamber-dependence of BPS spectrum.
% The shadow obstruction tower plays the role of the completion:
% it provides a chamber-independent, algebraic substitute for the
% mock modular counting function. The infinite tower of the BKM
% algebra is the chiral-algebraic shadow of the infinite Borcherds product.
% Fiber-to-global shadow depth jump: class G (tree level) -> class M
% (full non-perturbative answer) is the passage from sigma_3 = 0
% to sigma_3 != 0.

\subsection{Categorical $S$-matrix $\to$ $\mathrm{Sp}_4(\Z)$ (45 tests)}
% SCAFFOLDED (test_categorical_s_matrix_e3.py: 45 tests).
% Construction: S_{V,W} = Tr_W(R_{V,W} . R_{W,V}) (double braiding trace).
% Charge 1: PROVED. CY inversion g(z)g(-z)=1 gives S = 1 (trivial).
% Charge 2 diagonal (Fock): PROVED.
%   S_{(2)}(u) = g(u+h_2)*g(h_2-u), S_{(1,1)}(u) = g(u+h_1)*g(h_1-u).
%   S-matrix symmetry S(u) = S(-u) (even function).
% E_3 factorization: CONDITIONAL (Zamolodchikov ansatz
%   S^{E_3}(u,v) = S(u)*S(v)*S(u-v), even in both variables).
% The pipeline: E_3 braiding -> categorical S-matrix -> Sp_4(Z) action
%   -> Phi_10 as partition function.
% The Sp_4(Z) arises because K3 x E has lattice Lambda^{3,2} with
%   Sp_4(Z)/{+/-I_4} ~ O(Lambda^{3,2})_+/{+/-I_5} (Lemma k3e-lattice).
% Phi_10 connection: CONJECTURAL (depends on CY-A_3).

\subsection{$K3 \times E$ via factorization homology}
% Conj \ref{conj:fact-hom-k3xe}.
%   int_{K3} F|_{K3} = chiral algebra on E (character ~ Phi_10).
%   int_E (int_{K3} F) = partition function = Phi_10(tau).
%   Intermediate: toroidal KM enveloping chiral algebra on E.
% Three pathways to A_{K3xE} (Rem \ref{rem:6d-adds-to-k3xe}):
%   CY-to-chiral Phi, toroidal presentation, 6d hol theory.
% Agreement itself conjectural. kappa-spectrum: {2,3,5,24}.

%% ====================================================================
%% PART III: THE FRONTIER
%% ====================================================================

\section{The Frontier (Genuinely New)}

\subsection{Zamolodchikov tetrahedron equation for $U_{q,t}$}
% The E_3 braiding on C^3 predicts a TETRAHEDRON equation (3-simplex)
% R_{123} R_{145} R_{246} R_{356} = R_{356} R_{246} R_{145} R_{123}
% as the E_3 analogue of the YBE (E_2, 2-simplex).
% For U_{q,t}: the tetrahedron R-matrix acts on triple tensor
% products of Fock spaces. Connects to Kapranov-Voevodsky 2-categories.
% STATUS: equation written, no computational verification yet.

\subsection{Spin-3 chiral coproduct}
% Spin-2 scaffolded (23 tests). Spin-3: next target.
% At spin 3: W_3 current, first genuinely nonlinear OPE contribution.
% The coproduct Delta_ch(W_3) involves structure constants of W_{1+inf}
% beyond the Heisenberg sector. Requires full Yangian R-matrix at charge 3.
% Predicted: K-matrix modification at spin 3 detects class L vs class C.

\subsection{Two-parameter $R$-matrix $R(u,v)$}
% The 6d E_3 structure predicts R(u,v) with TWO spectral parameters
% (from two transverse directions beyond the Wilson line).
% Conj \ref{conj:chiral-qg-c3}(ii). Satisfies parametric YBE in (u,v).
% Relation to Maulik-Okounkov: MO gives R(u) (one parameter, from C^2);
% the second parameter v comes from the third C^3 direction.
% At v=0: reduces to the Yangian R-matrix of Thm \ref{thm:yang-r-matrix-ybe}.

\subsection{$\tau_3(n)$ as $E_3$ bar complex dimensions}
% Heuristic observation from holomorphic_cs_chiral_engine.py:
% The E_3 bar complex B_{E_3} of W_{1+inf} at arity n has
% dim B_{E_3,n} related to the 3D divisor function tau_3(n).
% Analogy: E_1 bar ~ p(n) (partitions), E_2 bar ~ d(n) (divisors),
%          E_3 bar ~ tau_3(n) (3-fold divisor function).
% tau_3(n) = #{(a,b,c): abc=n}. Generates 1/zeta(s)^3 in Dirichlet series.
% STATUS: heuristic, observed through n<=20. No proof.

\subsection{K3 double current algebra quantization}
% The chiral algebra int_{K3} F|_{K3} on E (Conj \ref{conj:fact-hom-k3xe}(iii))
% is a ``double current algebra'': currents from the K3 lattice Lambda_{K3}
% (rank 24) with two loop parameters (tau from E, sigma from K3 moduli).
% Quantization: the Omega-background on the K3 directions provides
% the deformation parameter. At the classical limit: lattice VOA V_{Lambda}.
% At generic deformation: a vertex algebra on E whose character is
% the Fourier-Jacobi expansion of Phi_10.
% STATUS: definition not constructed. Conditional on CY-A_3.

%% ====================================================================
%% PART IV: THE BOTTLENECK
%% ====================================================================

\section{The Bottleneck: Chain-Level $S^3$-Framing for CY-A$_3$}

\subsection{The obstruction}
% CY-A_2 is PROVED: the S^2-framing of HH_*(C) for CY_2 gives E_2-chiral.
% CY-A_3 requires S^3-framing. Thm \ref{thm:e1-stabilization-cy}: at d=3,
% Obs_eff(3) = 0 (pi_3(BU)=0, pi_3(BSp)=0). The obstruction GROUP is trivial,
% but the chain-level construction is not done.

\subsection{Existing approaches and the MO bypass}
% (A) Direct: construct the S^3-framing on HH_*(C) for CY_3 C at the chain level.
%     Blocked by: no explicit model for the S^3-framing of a DG category.
% (B) CoHA route: use SVYZ CoHA (E_1) + Drinfeld center to produce E_2.
%     BLOCKED by center-hocolim non-commutation (Prop \ref{prop:center-hocolim}):
%     the E_2 braided category cannot be assembled chart-by-chart.
%     Obstruction: C^3 (0), conifold (2), local P^2 (nonzero), K3xE (massive).
% (C) MO STABLE ENVELOPE BYPASS (NEW):
%     The MO stable-envelope R-matrix on Hilb^n(X) constructs E_2 braiding
%     directly from equivariant geometry, bypassing the Drinfeld center
%     entirely. For toric CY_3: works (Thm \ref{thm:mo-chiral-rmatrix}).
%     For K3 x E: works via Hilb^n(K3) (degree-24 structure function).
%     Limitation: requires a torus action or Hilbert scheme description;
%     does not work for general compact CY_3 without symmetry.
% (D) Holomorphic CS route: 6d E_3 on C^3 projected to E_1 on curve.
%     BYPASSES S^3-framing: the chiral algebra comes from factorization
%     homology, not from Phi. But 6d framework conjectural for compact M.
% (E) Relative chiral algebra route: for K3 x E, exploit the elliptic
%     fibration pi: S -> P^1 directly to build A_{K3,rel} without CY-A_3.
% Net: MO (C) and relative (E) provide concrete E_2 access for K3 x E;
% the abstract S^3-framing (A) remains open for general CY_3.

%% ====================================================================
%% PART V: OPEN PROBLEMS
%% ====================================================================

\section{Open Problems (Ranked by Impact)}

\begin{enumerate}
\item \textbf{Chain-level $S^3$-framing for CY-A$_3$.}
  Construct the functor $\Phi$ at $d=3$. All of Vol III downstream of CY-A depends on this. Impact: unlocks CY-B, CY-D, and the entire $K3 \times E$ programme. The MO stable envelope and relative chiral algebra routes bypass this for specific targets but not in general.

\item \textbf{6d algebraic framework for compact manifolds.}
  Extend the CFG factorization homology from $\C^3$ to $K3 \times E$. Would provide an alternative to Problem 1. Requires non-Lagrangian algebraic QFT on compact complex 3-folds.

\item \textbf{$E_3$ Koszul duality (Conj~\ref{conj:e3-koszul-duality}).}
  Prove $B_{E_3}$ has three commuting differentials and coproducts, and that $\cF^!$ has inverted parameters. The $(1+t)^{3g}$ formula for classes L, C provides strong evidence.

\item \textbf{Zamolodchikov tetrahedron equation for $U_{q,t}$.}
  Construct the tetrahedron $R$-matrix and verify the 3-simplex equation computationally at charge 2. First genuine $E_3$ invariant beyond the YBE.

\item \textbf{Two-parameter $R$-matrix $R(u,v)$.}
  Construct and verify the parametric YBE for the $E_3$ $R$-matrix on Fock spaces. Relates MO stable envelopes to the quantum toroidal braiding.

\item \textbf{Spin-3 chiral coproduct.}
  Extend the spin-2 scaffold (23 tests) to spin 3. First test of nonlinear OPE in the coproduct. Discriminates shadow class L from class C.

\item \textbf{$\int_{K3} \cF$ computation.}
  Compute the perturbative factorization homology over K3 and match to the Fourier--Jacobi expansion of $\Phi_{10}$. The Borcherds resummation result confirms the all-loop answer; extend to the non-perturbative regime.

\item \textbf{Categorical $S$-matrix to $\Phi_{10}$.}
  Close the gap between the charge-2 categorical $S$-matrix (proved, 45 tests) and the full $\mathrm{Sp}_4(\Z)$ modularity. Identify $\Phi_{10}$ as the partition function of the $E_3$ braided category.

\item \textbf{Algebraic vs topological $E_3$ over $\Z$.}
  Determine whether the two $E_3$ structures (Deligne conjecture vs configuration space) agree integrally or only rationally. Torsion discrepancy would affect the quantum toroidal algebra at roots of unity.

\item \textbf{Quantum group realization CY-C.}
  Construct $C(\frakg,q)$ as a chiral quantum group. Currently \texttt{\textbackslash begin\{conjecture\}} only. Requires resolution of Problems 1 or 2 as input.
\end{enumerate}

\end{document}
